%%%%%%%%%%%%%%%%%%%%%%%%%%%%%%%%%%%%%%%%%
% Original author:
% Adrien Friggeri (adrien@friggeri.net)
% https://github.com/afriggeri/CV
%
% Modified:
% Pol Cuscó
% Bibliography works with bibtex 3.3+ and biber.
%
% License:
% CC BY-NC-SA 3.0 (http://creativecommons.org/licenses/by-nc-sa/3.0/)
%
% Important notes:
% This template needs to be compiled with XeLaTeX and the bibliography, if used,
% needs to be compiled with biber rather than bibtex.
%
%%%%%%%%%%%%%%%%%%%%%%%%%%%%%%%%%%%%%%%%%

\documentclass[]{friggeri-cv} % Add 'print' as an option into the
                              % square bracket to remove colors from
                              % this template for printing

\usepackage{fontawesome}

%\addbibresource{bibliography.bib} % Specify the bibliography file to include publications

\begin{document}
\newfontfamily{\FA}{FontAwesome}

% Your name and current job title/field
\header{pol}{cuscó}{bioinformatician \&\& machine learning geek}

%----------------------------------------------------------------------------------------
%	SIDEBAR SECTION
%----------------------------------------------------------------------------------------

\begin{aside} % In the aside, each new line forces a line break
  \section{contact}
  %{\FA \faUser}
  %Pol Cuscó Pons
  %{\FA \faEnvelope}
  \includegraphics[scale=0.03]{emaillogo.jpeg}
  \href{mailto:polcusco@gmail.com}{polcusco@gmail}
  %{\FA \faLinkedinSign}
  \includegraphics[scale=0.02]{linkedinlogo.jpeg}
  \href{http://www.linkedin.com/in/pol-cuscó-80246963}{linkedin/pol-cusco}
  \includegraphics[scale=0.07]{researchgatelogo.png}
  \href{https://www.researchgate.net/profile/Pol\_Cusco2}{researchgate/Pol\_Cusco2}
  %{\FA \faGithubSign}
  \includegraphics[scale=0.08]{githublogo.jpeg}
  \href{https://github.com/nanakiksc}{github.com/nanakiksc}
  \includegraphics[scale=0.6]{phonelogo.png}
  %{\FA \faPhone}
  +34 684 200 939

  \section{languages}
  catalan \& spanish
  english
  portuguese
  \section{programming}
  {\color{red} \FA \faHeart} R
  Python
  C \& C\#
  HTML \& CSS
  Shell scripting
  \LaTeX
\end{aside}

%----------------------------------------------------------------------------------------
%	INTERESTS SECTION
%----------------------------------------------------------------------------------------

\section{interests}

genomics, machine learning, big data, data visualization

%----------------------------------------------------------------------------------------
%	EDUCATION SECTION
%----------------------------------------------------------------------------------------

\section{education}

\begin{entrylist}

\entry
{since 2013}
{PhD {\normalfont in Biomedicine}}
{Universitat Pompeu Fabra}
{\emph{Machine learning approach to the study of chromatin}\\
\href{https://www.tdx.cat/handle/10803/565685}{{\FA \faExternalLink} Thesis.}}

\entry
{2010--2011}
{MSc {\normalfont in Pharmaceutical Industry \& Biotechnology}}
{Universitat Pompeu Fabra}
{}

\entry
{2005--2010}
{BSc {\normalfont in Biology}}
{Universitat Pompeu Fabra}
{Specialization in Human Biology}


\entry
{2012--2013}
{Research Scholar}
{Massachusetts Institute of Technology, Cambridge (MA)}
{I was enrolled at the Department of Mechanical Engineering,
  where I conducted research on autonomous underwater vehicles and
  underwater wireless communications.}

\end{entrylist}

%----------------------------------------------------------------------------------------
%	WORK EXPERIENCE SECTION
%----------------------------------------------------------------------------------------

\section{experience}

\begin{entrylist}

  \entry
    {2019}
    {Gastrointestinal and Endocrine Tumor Group}
    {Vall d'Hebron Institute of Oncology}
    {\emph{Bioinformatician}\\
    I developed and maintained NGS analysis pipelines,
    working mostly with whole-exome sequencing data from both
    tumors and cfDNA. I also wrote scripts for custom
    analyses and performed data visualization tasks.
    }
  \entry
    {2017--2019}
    {Bioinformatics Unit}
    {Vall d'Hebron Institute of Oncology}
    {\emph{Bioinformatician}\\
    I also worked with different research groups in the institute
    and gave them bioinformatics support. I helped them decide
    which data analyses they needed in their projects and then
    I conducted them. Is was also my responsibility to
    communicate the results of such analyses to help guide
    their research.
    }
  \entry
    {2013--2017}
    {Genome Architecture Lab}
    {Centre for Genomic Regulation}
    {\emph{PhD student}\\
    I worked from raw genomics data, especially ChIP-seq.
    I used both supervised
    and unsupervised machine learning methods and developed of our
    own (see publications). I have a solid grasp of best practices in
    machine learning and how to avoid common pitfalls (see
    publications).\\
    \href{http://www.genomearchitecture.com}{{\FA \faExternalLink} Lab website.}}
\end{entrylist}
\begin{entrylist}
  \entry
    {2011--2012}
    {Neuropharmacology Lab}
    {Department of Experimental and Health Sciences, UPF}
    {Technician}
  \entry
    {2011}
    {Central Nervous System Unit}
    {Pharmaceutical R\&D Centre, Ferrer Internacional}
    {Researcher during my MSc internship}
  \entry
    {2010}
    {Keypoint CRO}
    {Grupo Keypoint}
    {Clinical Research Associate during my BSc internship}
  \entry
    {2008}
    {Neuropharmacology Lab}
    {Department of Experimental and Health Sciences, UPF}
    {Student during a summer internship}

\end{entrylist}

%----------------------------------------------------------------------------------------
%	PUBLICATIONS SECTION
%----------------------------------------------------------------------------------------

\section{publications}


\begin{itemize}
    \item Cuscó P, Filion GJ. Zerone: a ChIP-seq discretizer for
        multiple replicates with built-in quality control.
        Bioinformatics. 2016;32(19):2896-902.
    \item Corrales M, Cuscó P, Usmanova DR, Chen HC, Bogatyreva NS,
        Filion GJ, Ivankov DN. Machine Learning: How Much Does It
        Tell about Protein Folding Rates? PLoS One.
        2015;10(11):e0143166.
    \item Zorita E, Cuscó P, Filion GJ. Starcode: sequence clustering
        based on all-pairs search. Bioinformatics. 2015;31(12):1913-9.
\end{itemize}

%----------------------------------------------------------------------------------------
%	COURSES SECTION
%----------------------------------------------------------------------------------------

\section{courses}

\begin{entrylist}
  \entry
    {2017}
    {From Science to Market}
    {Universitat de Barcelona / Universitat Politècnica de Catalunya}
    {A course on entrepreneurship focused on technology-based
    companies.}
  \entry
    {2014}
    {Algorithms: Design and Analysis, Part 1}
    {Stanford University}
    {An in-depth review of classical algorithms and data structures.}
  \entry
    {2013}
    {Machine Learning}
    {Stanford University}
    {A reference on the fundamentals and good practices of
    machine learning.}
  \entry
    {2012}
    {Web Fundamentals, JavaScript, jQuery, PHP, Python}
    {Codecademy}
    {Several courses on different programming languages.}
  \entry
    {2012}
    {Introduction to Computing Principles}
    {Stanford University}
    {An introduction to the key ideas of computing.}
\end{entrylist}

%----------------------------------------------------------------------------------------

\end{document}
